% Chapter: Conclusions and Future Work
% Pathway A framing: characterisation of a structured BZ reaction-diffusion
% medium as a trainable analog computing substrate.

\section{Summary of Contributions}
\label{sec:conc_summary}

The aim of this thesis was to characterise how a structured reaction--diffusion excitable medium processes information, treating the medium as an analog computing substrate whose structure can be understood and ultimately designed for computation.
The work makes five main contributions.

First, it justifies the choice of substrate.
A systematic comparison of candidate physical computing families---acoustic, thermal, well-mixed chemical, DNA strand-displacement, and reaction--diffusion excitable media---concludes that a structured Belousov--Zhabotinsky (BZ) medium offers the best combination of native nonlinearity, continuous parallel space, a compact validated model, and experimental accessibility (\refChap{chap:substrates}).

Second, it implements and verifies a numerical solver for the two-variable Tyson--Fife Oregonator on a two-dimensional Cartesian grid, extended with no-flux walls, a spatial light-suppression field $\phi(x,y)$, and an anisotropic activator diffusion tensor.
The solver is verified by the method of manufactured solutions and by grid and timestep convergence studies; the Oregonator pulse speed is grid-converged within a $2.8\%$ band at the working resolution (\refSec{sec:res_verification}).

Third, it establishes a black-box characterisation protocol.
The substrate is stimulated at defined input ports and read at fixed probes; every reported result is an input--output transfer function, with no inspection of the internal field used for the measurement (\refSec{sec:meth_protocol}).

Fourth, it measures a set of native transfer functions that together form a component datasheet for the substrate: near-ideal pulse transmission in straight channels, a refractory-limited maximum 1:1 following rate of $0.282$ per time unit, an inhibition gate A~AND~(NOT~B) with separation ratios above $200\times$, and continuously tunable routing delays via anisotropic diffusion patches (\refChap{chap:results}).

Fifth, it demonstrates that gradient-free optimisation can discover a working routing field.
A fixed-wall T-junction with an optimised illumination field $\phi(x,y)$ routes a single input pulse to one output probe while keeping the opposite probe below threshold (\refSec{sec:res_tjunc_train}).
The trained field recovers the same high-$\phi$ blocking-barrier physics that hand-designed fields used earlier, so the result is better read as a feasibility demonstration than as evidence of complex learned computation.

\section{Limitations}
\label{sec:conc_limitations}

The results are subject to several limitations that should be stated openly.
The geometry is idealised to two dimensions; real BZ layers are thin but finite, and three-dimensional scroll-wave dynamics or boundary effects are not captured.
The Oregonator is a reduced model of the full Field--K\"or\"os--Noyes chemistry; it reproduces the excitable phenomenology but does not resolve all chemical species or recipe-dependent kinetic details.
The solver is explicit on a regular grid, so absolute speeds and timings carry a spatial discretisation error that is quantified but not eliminated at the working resolution.
The physical speed of chemical waves---of order $0.1$~mm/s---is slow compared with electronic logic; the substrate's advantage is spatial parallelism, not latency.
All results are numerical; no wet-lab validation in a photosensitive BZ layer has been performed.
Finally, the work stops at characterisation: there is no closed optimisation loop that automatically designs a wall geometry or $\phi(x,y)$ for an arbitrary target input--output map.

\section{Future Work}
\label{sec:conc_future}

The most direct next step is to close the loop from characterisation to training.
Gradient-free or surrogate-based optimisers should be used to design wall geometry, the light field $\phi(x,y)$, and the anisotropic diffusion tensor field for target input--output maps, using the transfer functions measured here as the feasibility constraints on the design space.
Because the forward solver is expensive and the loss landscape is noisy and non-differentiable, Bayesian optimisation or hybrid model-based methods are natural candidates.

A second direction is uncertainty quantification and robust design.
The parameters of the Oregonator, the illumination field, and the initial conditions are all uncertain in a real experiment.
Polynomial chaos expansion methods have already been applied to molecular signals in fluid-dynamic channels \cite{abbaszadeh2019uncertainty}, and their extension to the reaction--diffusion channel would let trained fields be validated against parameter uncertainty before any wet-lab test \cite{guo2021molecular}.

A third direction is physics-informed inference from sparse probes.
Rather than treating the internal field as hidden, a training or calibration stage could infer the effective medium parameters from probe time series, turning the substrate into a physically informed communication channel whose state can be tracked from limited observations \cite{lin2023diversity}.
Related to this, Koopman reduced-order models and sparse probe placement could compress the high-dimensional field dynamics into a low-dimensional state representation suitable for real-time readout and control.

Longer term, experimental validation in a photosensitive ruthenium-catalysed BZ layer would test whether the measured transfer functions survive the recipe dependence, optical projection resolution, and thermal drift of a real chemical implementation.

\section{Concluding Remarks}
\label{sec:conc_remarks}

This thesis began from the premise that a fluid or chemical medium can compute continuously and in parallel, and that programming such a medium means shaping its structure and property fields.
The central argument is that characterisation must precede training: before optimising a structured reaction--diffusion substrate for a task, one must know what it already does.
The four transfer functions reported here provide that baseline.
They show that a BZ-type medium is not a general-purpose digital computer, but a slow, parallel, threshold-logic substrate with well-defined wiring, bandwidth, logic, and routing primitives.
Those primitives are exactly the levers from which trainable chemical processors can be built, and the characterisation-first approach offers a general recipe for turning physical substrates into programmable analog machines.
